\subsection{Independent checks of the full-group arguments}

\statement{C.5}{Proved}{Primitive-homology centralizers and the original torus cusp}
If $F$ is free profinite and $a\in F$ admits a continuous
character $\rho:F\to\widehat{\mathbb Z}$ with $\rho(a)=1$,
then, for every nonzero integer $n$,
\[
 C_F(a^n)=N_F(\overline{\langle a\rangle})
       =\overline{\langle a\rangle}\cong\widehat{\mathbb Z}.
                                                               \tag{C.5.1}
\]
For $F=\widehat F(x,y)$ and $c=[x,y]$, one also has
$C_F(c)=N_F(\overline{\langle c\rangle})=\overline{\langle c\rangle}$.
The first assertion requires primitive homology, not membership
in a free basis.

\begin{proof}
The character makes $t\mapsto a^t$ injective.  Suppose that
$k\ne1$ lies in $C_F(a^n)\cap\ker\rho$.  Some primary
component of the procyclic group generated by $k$ is
nontrivial, hence is $\mathbb Z_p$, since free profinite
groups are torsion-free.  It commutes with the primary
component $\mathbb Z_p$ of $\langle a^n\rangle$.  Their
intersection is trivial: $\rho$ vanishes on one and is
injective on the other.  Their product embeds the closed
group $\mathbb Z_p^2$ in $F$, contradicting
$\operatorname{cd}_pF=1$.  Thus $\rho$ is injective on
the centralizer.  For any $h$ in it,
$ha^{-\rho(h)}$ belongs to its kernel, proving the
centralizer formula.  A normalizer sends $a$ to $a^u$;
applying $\rho$ gives $u=1$, proving the normalizer formula.

For the cusp use $U=\ker(F\to C_2)$, with $x\mapsto1$,
$y\mapsto0$.  Its Schreier basis is $x^2,y,xyx^{-1}$,
and $c=(xyx^{-1})y^{-1}$ is a member of a free basis of
$U$.  Thus $C_U(c)=\overline{\langle c\rangle}$.  If
$h\in C_F(c)\setminus U$, then $h^2=c^t$ for some
$t\in\widehat{\mathbb Z}$.  Abelianization gives $2[h]=0$,
so $[h]=0$, contradicting $h\notin U$.  This proves
the centralizer assertion.  A normalizer sends $c$ to
$c^u$.  In the upper unitriangular Heisenberg quotient
over $\widehat{\mathbb Z}$, the image of $c$ is the
central coordinate one; conjugation leaves it unchanged.
Thus $u=1$, proving the remaining assertion.
\end{proof}

\statement{C.6}{Proved}{The finite-coefficient duality behind LM4 D.74--D.78}
Let $(F;L_1,\ldots,L_m)$ be the full completion of an
oriented compact surface with nonempty boundary and
negative Euler characteristic.  Let $E=F\rtimes C_2$
preserve each peripheral conjugacy class with exponent
$-1$, and suppose $D_j=N_E(L_j)=L_j\rtimes_{-1}C_2$.
Put $k=\mathbb F_2$, let $u:E\to k$ be the quotient
character, and define
\[
 \Delta=\ker\left(\bigoplus_j k[[E/D_j]]\longrightarrow k\right),
 \qquad \mathscr H_M=\operatorname{Hom}_{k,\mathrm{cont}}(\Delta,M).
\]
For every finite continuous $E$-module $M$, absolute
cohomology with coefficient $M$ and relative cohomology
with coefficient $M^\vee$ have a natural perfect pairing
after localization by $u$.  For self-dual $M$, the
localized boundary restriction has half the target rank:
\[
 \operatorname{rank}\left(
 H^\bullet(E,M)[u^{-1}]_0\longrightarrow
 \bigoplus_jH^\bullet(D_j,M)[u^{-1}]_0\right)
 =\frac12\sum_j\dim H^\bullet(D_j,M)[u^{-1}]_0.               \tag{C.6.1}
\]
The subscript zero denotes degree zero after localization.
The relative groups here use the original peripheral
normalizers.

\begin{proof}
The map $F/L_j\to E/D_j$, $fL_j\mapsto fD_j$, is an
$F$-equivariant homeomorphism: $E=FD_j$ and
$F\cap D_j=L_j$.  Restriction of $\Delta$ is therefore
the original surface-pair augmentation module.  For a
finite coefficient $M$, put $U=H^0(F,M)$, $V=H^1(F,M)$.
Surface-pair duality gives the natural $C_2$-equivariant
identifications
\[
 H^0(F,\mathscr H_{M^\vee})=V^\vee,
 \qquad H^1(F,\mathscr H_{M^\vee})=U^\vee.                  \tag{C.6.2}
\]
The orientation character is trivial modulo two.  The
full-completion comparison and cup-product duality are
Propositions~6.16 and 5.7 of the accepted version of
\cite{CEWilkes2019}.  Relative cohomology in degree
$n\geq1$ is $H^{n-1}(E,\mathscr H_M)$; its long exact
sequence contains the actual restrictions to the $D_j$.

The two ordinary Hochschild--Serre spectral sequences
for $M$ and $\mathscr H_{M^\vee}$ have fibre rows
$(U,V)$ and $(V^\vee,U^\vee)$.  In each positive
quotient column, cohomology of a finite $C_2$-module $W$
is
\[
 \overline W=\ker(1+\sigma)/\operatorname{im}(1+\sigma).
\]
Evaluation gives a perfect pairing
$\overline W\times\overline{W^\vee}\to k$.
The only possible differentials are
$d:\overline V\to\overline U$ and
$d':\overline{U^\vee}\to\overline{V^\vee}$.
Coefficient evaluation gives
$M\otimes\mathscr H_{M^\vee}\to\mathscr H_k$.
The product of the two fibre-row-one classes lands in
$H^2(F,\mathscr H_k)=0$.  The derivation identity for
$d_2$, followed by fundamental-class evaluation, gives
\[
                  \langle dv,u^\vee\rangle
                    =\langle v,d'u^\vee\rangle.              \tag{C.6.3}
\]
Thus the differentials are transposes.  No vanishing
of either differential for nontrivial coefficients
has been assumed.

The surviving rows pair perfectly as
$\operatorname{coker}d$ with $\ker d'$ and $\ker d$
with $\operatorname{coker}d'$.  For trivial coefficients
the section makes inflation injective, so the absolute
differential and its transpose vanish.  The top relative
fibre class supplies compatible traces
$H^n(E,\{D_j\};k)\to ku^{n-2}\to k$ in high degrees.
Cup product followed by this trace pairs the two
filtered spaces above.  The two filtration subspaces
pair to zero by filtration degree; the opposite
subspace--quotient pairings are precisely the two
perfect pairings already obtained.  A nonzero vector
in a filtration subspace pairs with a lift of the
opposite quotient, and a vector with nonzero quotient
pairs with the opposite subspace.  This proves
nondegeneracy on the full spaces.  Multiplication by
$u$ is an isomorphism in sufficiently high degree,
as is seen on the positive-column cyclic complex.
It gives the stated localized perfect pairing.

For self-dual $M$, denote the localized absolute,
boundary and relative spaces by $\mathcal A,
\mathcal B,\mathcal R$.  Their cyclic exact sequence
is $\mathcal A\to\mathcal B\to\mathcal R\to\mathcal A$,
with the connecting map divided by $u$.  The dimensions
of $\mathcal A$ and $\mathcal R$ agree by duality.
If the three ranks are $x,y,z$, exactness gives
$\dim\mathcal A=x+z$, $\dim\mathcal B=x+y$ and
$\dim\mathcal R=y+z$.  Hence $x=y$, proving (C.6.1).
The pairing is natural and balanced over the localized
scalar ring by the cup-product and connecting-map
compatibility in Section~3 of \cite{CEWilkes2019}.
\end{proof}

\statement{C.7}{Proved}{The full pants section centralizers in LM4 D.74}
For the normalized involution and three section classes
of R1, every section has trivial centralizer in the full
$F=\widehat F(x,y)$.  Thus $N_E(S_i)=S_i$ for all
three classes.  The equation $f=h^{-1}\alpha(h)$ has
exactly one solution in $F$, and that solution lies in
$F'$.  The extension $E=F\rtimes_\alpha C_2$ is real
projective.  If $\alpha\in\widehat{GT}_0$, the unique
solution also satisfies $h\theta(h)=1$, where $\theta$
exchanges $x,y$.

\begin{proof}
The pants peripheral generators have primitive homology,
so C.5 identifies each geometric normalizer with the
procyclic cusp itself.  A cusp-normalizer lift in $E$
inverts that subgroup.  Its square lies in the cusp and
is fixed by inversion, hence is one.  The three
peripheral normalizers are therefore those of C.6.
Each has two local section classes with self-normalizers.
R1 supplies their triangle incidence: each of the three
global section classes occurs at exactly two endpoints.

Write $C_i=C_F(S_i)$; its full normalizer is
$C_i\times S_i$, with factors separated by the
quotient label.  Theorem~2 of
\cite{CEKochloukovaZalesskii2007} applies to $E$ and
to its peripheral normalizers: their open geometric
groups have cohomological dimension one and no
$C_2^2$ embeds, since every finite subgroup injects
into $C_2$.  It gives a dense restriction embedding
in degrees above one.  The source is finite by the
two-row spectral sequence, so this is an isomorphism.
After localization, restriction to the six local
factors is the direct sum of three maps
\[
 r_i:H^\bullet(C_i\times S_i,M)[u^{-1}]_0
       \longrightarrow H^\bullet(S_i,M)[u^{-1}]_0,
\]
each occurring twice, with its actual conjugating
coefficient transport.  C.6 says the total rank is
the sum of the three target dimensions.  Each rank
is at most its target dimension; hence every $r_i$
is onto for every finite self-dual permutation module.

If $1\ne h\in C_i$, choose a finite quotient $D$ of
$E$ in which $h$ is not in the image of $S_i$, and
take $M=k[D/\overline S_i]$.  In positive cyclic
cohomology the $S_i$-fixed cosets give a basis, while
the two-element orbits contribute zero.  The element
$h$ moves the identity-coset basis vector to a
distinct fixed-coset vector.  However, the image of
restriction from $C_i\times S_i$ is $C_i$-invariant,
by naturality under inner conjugation.  This
contradicts surjectivity of $r_i$, proving $C_i=1$.

R1 gives a solution $h\in F'$.  If $h'$ is another,
the equations imply $\alpha(h'h^{-1})=h'h^{-1}$;
the centralizer just computed gives $h'=h$.
For $\widehat{GT}_0$, use
$\theta(f)=f^{-1}$ and
$\theta\alpha\theta=\operatorname{Inn}(f)\alpha$.
Applying them to the coboundary equation shows that
$\theta(h)^{-1}$ is another solution; uniqueness
proves $h\theta(h)=1$.
The set of involutions is closed because it lies
in the open nonidentity quotient coset.  The open
subgroup $F$ is projective and every involution is
self-centralizing.  Proposition~2.2 of
\cite{CEHaran1993} therefore proves real projectivity:
finite embedding problems whose involution images
admit involution lifts have continuous weak solutions.
This does not assert a free-product presentation.
\end{proof}

\statement{C.8}{Proved}{Negative torus section counts, cusp fusion and the centralizer in LM4 D.75}
Let an outer involution of $F=\widehat F(x,y)$ preserve
the commutator cusp $c=[x,y]$ with full multiplier $-1$.
Its centre-free extension has a cusp-normalized splitting
\[
 E=F\rtimes\langle v\rangle,
 \qquad v^2=1,\qquad vcv^{-1}=c^{-1}.
\]
Write $A_2$ for its action on $H_1(F,\mathbb Z_2)$.
If $A_2\not\equiv I\pmod2$ there is one section class
and its geometric centralizer is trivial.  If
$A_2\equiv I\pmod2$ there are two classes: the cusp
class has trivial geometric centralizer and the other
has $\dim_{\mathbb F_2}H^1(C_F(w),\mathbb F_2)=1$.
The two local cusp reflection classes are globally
fused.  In particular there is a unique $g\in F$ with
$gvg^{-1}=cv$.

\begin{proof}
C.5 gives the geometric cusp normalizer
$\overline{\langle c\rangle}$.  Any lift $v$ in
the outer-extension cusp normalizer inverts $c$.
Its square is $c^t$ and is fixed by inversion,
whence $2t=0$ and $v^2=1$.  This proves the stated
splitting and the full dihedral cusp normalizer.
The Heisenberg quotient in C.5 shows
$\det A_2=-1$.  Every finite subgroup of $E$ injects
into $C_2$, and each section $S$ has normalizer
$C_F(S)\times S$.

Put $k=\mathbb F_2$.  If $A_2\equiv I\pmod2$,
the operator $(I-A_2)/2$ is an integral idempotent.
Its image and kernel have rank one, so the integral
character module has a basis $\varphi_+,\varphi_-$
with actions $+1,-1$.  Their reductions extend to
characters $a,b:E\to k$, both zero on $v$.
The two-row spectral sequence collapses and, for
$n\geq1$, has basis
\[
                  H^n(E,k)=\langle u^n,u^{n-1}a,u^{n-1}b\rangle.
\]
The first character lifts to $\mathbb Z/4$, so
$a^2=0$.  The cochain
$\widetilde b(hv^\epsilon)=\varphi_-(h)\pmod4$
has coboundary $2\epsilon\varphi_-(h')$ on
$(hv^\epsilon,h'v^\eta)$.  Its Bockstein gives
$b^2=ub$.  Write $ab=\epsilon ua+\delta ub+
\gamma u^2$.  Restriction to $v$ gives $\gamma=0$;
multiplication by $a$ and independence in degree
three give $\delta=0$.  Hence localization gives
\[
 H^\bullet(E,k)[u^{-1}]_0
 \cong k[A,B]/(A^2,B^2-B,AB-\epsilon A)
 \cong k[\nu]/(\nu^2)\times k.                              \tag{C.8.1}
\]
The last isomorphism takes $A$ to $(\nu,0)$ and
$B$ to $(\epsilon,1-\epsilon)$.
If $A_2\not\equiv I$, its dual mod-two module is
the regular $k[C_2]$-module, whose positive
cohomology vanishes.  The same two-row calculation
then gives $H^n(E,k)=ku^n$ for $n\geq2$ and
localized ring $k$.

The finite-coefficient restriction theorem used in
C.7 identifies these localized rings with
\[
 \prod_{[S]}\bigl(k\oplus H^1(C_F(S),k)u^{-1}\bigr).        \tag{C.8.2}
\]
Each second summand has square zero, since a closed
subgroup of $F$ has cohomological dimension at most
one.  The local factors and their radicals in
(C.8.1)--(C.8.2) give the asserted section counts
and the multiset of first-cohomology dimensions
$\{1,0\}$ or $\{0\}$.  In the two-class case the
values of $b$ distinguish the classes, as seen
on the two residue fields in (C.8.1).  The local
reflections $v,cv$ have the same value under every
abelian character because $c\in[F,F]$.  They are
therefore globally conjugate in either case.
Conjugacy by $E$ is the same as conjugacy by $F$
for sections: multiply a conjugator by the source
reflection to correct its quotient label.

Apply C.6 to the one actual torus cusp.  Its two
local factors are the same global section factor,
with the transport supplied by a chosen conjugator
from $v$ to $cv$.  The other global factor, when
present, has zero boundary restriction.  The
half-rank identity says precisely that
\[
 H^\bullet(C_F(v)\times\langle v\rangle,M)[u^{-1}]_0
        \longrightarrow H^\bullet(\langle v\rangle,M)[u^{-1}]_0
\]
is onto for every self-dual permutation coefficient.
The fixed-coset witness from C.7 forces $C_F(v)=1$.
This also identifies which factor in (C.8.1) is
the cusp class.  All cusp reflections are in that
class, so all have trivial geometric centralizer.
Two conjugators from $v$ to $cv$ differ by
$C_F(v)$, proving uniqueness of $g$.  In the
one-class case all involutions are self-centralizing,
so Haran's criterion also makes $E$ real projective.
\end{proof}

\statement{C.9}{Proved}{The specified section augmentation module in LM4 D.76}
In the pants extension of C.7 choose one actual section
$T_i=\langle t_i\rangle$ from each class, and put
$R=\mathbb F_2[[E]]$.  Then the compact left module
\[
 \Delta_T=\ker\left(\bigoplus_{i=0}^2
                 \mathbb F_2[[E/T_i]]\longrightarrow\mathbb F_2\right)
\]
is isomorphic to $R^2$.  For each finite $E$-module $M$,
$\dim\operatorname{Hom}_E(\Delta_T,M)=2\dim M$.
For the specified identity-coset vectors $e_i$ and
$H=\overline{\langle T_0,T_1,T_2\rangle}$, the distinct
incidence map satisfies
\[
 \begin{split}
 d_T:R^2&\longrightarrow\Delta_T,
   \quad(r_1,r_2)\longmapsto r_1(e_1-e_0)+r_2(e_2-e_0),\\
 \operatorname{coker}d_T&\cong
       \ker\bigl(\mathbb F_2[[E/H]]\longrightarrow\mathbb F_2\bigr).
 \end{split}                                                  \tag{C.9.1}
\]
Thus $d_T$ is invertible exactly when the chosen sections
generate $E$; freeness of $\Delta_T$ alone does not imply
that conclusion.

\begin{proof}
Each $F\to E/T_i$ is an $F$-equivariant homeomorphism.
On restriction to $R_F=k[[F]]$ the augmentation kernel
is $I_F\oplus R_F^2\cong R_F^4$, since
$I_F=R_F(x-1)\oplus R_F(y-1)$.  The latter identity
can be checked by applying continuous Hom to any
finite module: a crossed homomorphism from the full
free group is uniquely specified by arbitrary values
on $x,y$.  Finite quotients separate maps.  In
particular $\Delta_T$ is finitely generated over $R$.

The exact coefficient sequence
\[
 0\to M\to M\otimes_k k[E/F]\to M\to0
\]
has connecting map cup product by $u$.  Its middle
term is induced from $F$, so Shapiro and
$H^2(F,M)=0$ show that $u:H^1(E,M)\to H^2(E,M)$
is onto.  C.7 and the restriction theorem identify
$H^2(E,M)$ with $\bigoplus_i H^2(T_i,M)$.  Cup
product by $u$ is an isomorphism from $H^1(T_i,M)$
to $H^2(T_i,M)$ by the cyclic resolution.  Consequently
the degree-one restriction is onto as well.  The
augmentation long exact sequence now gives
$\operatorname{Ext}^1_R(\Delta_T,M)=0$ for every
finite $M$.

This implies projectivity in the compact module
category.  To see the finite-limit point, choose
$R^d\twoheadrightarrow\Delta_T$.  For each finite
quotient of this epimorphism the pulled-back
extension over $\Delta_T$ has finite kernel and
splits by that Ext vanishing.  The finite sets of
lift maps have the finite-intersection property,
by passing to a common quotient refinement.
Compactness supplies a compatible splitting of
the original epimorphism.

Write $e_j=\dim H^j(E,M)$, $f_j=\dim H^j(F,M)$,
and $s_i=\dim M^{T_i}$.  The coefficient sequence
and Shapiro give
$2e_1-e_2=2e_0-f_0+f_1=2e_0+\dim M$.
The cyclic resolution gives
$\dim H^1(T_i,M)=\dim H^2(T_i,M)=2s_i-\dim M$.
Brown restriction gives $e_2=2\sum_i s_i-3\dim M$.
In the augmentation Hom sequence the degree-one
restriction is onto; substituting these identities
gives exactly
$\dim\operatorname{Hom}_E(\Delta_T,M)=2\dim M$.

For every open normal $N\lhd E$, base change gives
a finite projective module over $R_N=k[E/N]$.
Its Hom multiplicities into every simple $R_N$-module
are those of $R_N^2$.  The tops, and hence their
projective covers, agree; thus this module is
$R_N^2$.  Compatible choices of isomorphisms exist
by compactness: any finitely many compatibility
conditions are satisfied by reduction of an
isomorphism at a common finer quotient.  Their
inverse limit proves $\Delta_T\cong R^2$.

Finally quotienting the direct sum of permutation
modules by the image of $d_T$ identifies all three
identity-coset vectors.  The resulting vector is
fixed by exactly the generated subgroup $H$, giving
$k[[E/H]]$ and the augmentation-kernel identity
(C.9.1).  This kernel vanishes exactly when $H=E$.
If $H=E$, $d_T$ is a surjective endomorphism of
$R^2$ after the constructed identification; its
reductions at every finite group ring are
bijective, so it is an isomorphism.  This proves
the precise incidence criterion.
\end{proof}

\statement{C.10}{Proved}{The proper three-section example of LM4 D.77}
In the full profinite free product
$E_0=C_2(a)\amalg C_2(b)\amalg C_2(d)$ put
\[
 x=ba,\quad y=da,\quad \delta=\operatorname{Ad}(a),\quad
 g=x[x,y]x^{-1},\quad Y=g y\delta(g)^{-1}.
\]
The orientation kernel is $F_0=\widehat F(x,y)$.
The sections $a,b,gdg^{-1}$ represent its three
section classes, yet generate a proper closed subgroup
$H<E_0$.  Its orientation kernel is
$K=\overline{\langle x,Y\rangle}<F_0$, although its
image in every $F_0^{(p)}$ is full.  The incidence
cokernel of C.9 has a five-dimensional finite quotient.
The element $Y$ is not conjugate to any unit power of $y$.

\begin{proof}
The inverse maps $(a,b,d)\mapsto(a,xa,ya)$ and
$(x,y,a)\mapsto(ba,da,a)$ identify $E_0$ with
$\widehat F(x,y)\rtimes_\delta C_2$, where
$\delta(x)=x^{-1}$ and $\delta(y)=y^{-1}$.
The three factors have distinct images in $C_2^3$,
so conjugating the third retains its section label.
The identities $gdg^{-1}=Ya$ and $\delta(Y)=Y^{-1}$
show $H=K\rtimes\langle a\rangle$.  Since $g\in F_0'$,
the pair $x,Y$ induces the identity on abelianization.
It therefore generates every elementary abelian
Frattini quotient of the free pro-$p$ group.  A
proper closed subgroup would lie in an index-$p$
maximal subgroup in some finite $p$-group quotient,
and be annihilated by a nonzero character to $C_p$.
This is impossible, proving the all-prime assertion.

Use permutation composition $(uv)(i)=u(v(i))$ and
the continuous finite quotient
\[
 a\mapsto(12)(34),\qquad b\mapsto(13)(25),\qquad
 d\mapsto(14)(23).
\]
Multiplication gives
\[
 x\mapsto(1\,5\,2\,3\,4),\quad y\mapsto(13)(24),\quad
 g\mapsto(1\,3\,2),\quad gdg^{-1}\mapsto a,\quad Y\mapsto1.
\]
The original triple and $x,y$ each generate $A_5$;
the chosen sections generate $D_{10}$ and $x,Y$
generate $C_5$.  This proves both proper inclusions.
The coset map $E_0/H\to A_5/D_{10}$ gives a quotient
of the incidence cokernel of dimension $6-1=5$.
Finally $y$ and all its unit powers have image of
order two, whereas $Y$ has image one, precluding
their conjugacy.  The script
\texttt{code/check\_D77.py} independently enumerates
all these generated subgroups and verifies the
permutations.  The Frattini argument, rather than
a finite list of prime tests, proves the primary claim.
\end{proof}

\statement{C.11}{Proved}{The original negative transporter is homologically primitive}
In C.8 the transporter $g$ has primitive abelianization
and generates the full fixed line of $F^{\mathrm{ab}}$.
In the two-class case put $B=C_F(w)$ for the other
section.  Its homology image is exactly that line.
There is a continuous epimorphism
$\chi:E\to\widehat{\mathbb Z}$ with $\chi(g)=1$ and
$\chi(B)=\widehat{\mathbb Z}$; one can choose
$a\in B$ with $[a]=[g]$.  Killing all torsion in
$E$ gives precisely this torsion-free quotient.
In the one-class case the torsion-free quotient
is also $\widehat{\mathbb Z}$, but its primitive
character takes $g$ to $2$.

\begin{proof}
The identity $c=g\alpha(g)^{-1}$ gives
$[g]\in M^\alpha$, where $M=F^{\mathrm{ab}}$.
At each prime the negative involution has a
rank-one fixed line.  For an arbitrary odd integer
$m\geq3$ choose a surjective invariant character
$\eta:F\to C_m$ and extend it by $\eta(v)=0$.
The odd-primary $+1,-1$ eigenspace decomposition
provides this character.  Put $F_m=\ker\eta$ and
$E_m=F_m\rtimes\langle v\rangle$.  Its genuine
surface cover has type $(1,m)$: the cusp has
zero image, so all $m$ cusps have degree one,
and Euler characteristic is multiplied by $m$.

Deck averaging over the odd group $C_m$ splits
$H_1(F_m,\mathbb Z_2)$ into the rank-two compact
torus part and the rank-$(m-1)$ boundary
augmentation part.  The latter consists of the
$m$ boundary loops with their sole sum relation.
The involution fixes each boundary label and
acts by $-1$ there.  Modulo two every boundary
character is invariant and extends to $E_m$ by
assigning zero to $v$.

The cusp sections split into $m$ classes in
$E_m$, since C.8 gives $C_F(v)=1$.
The two endpoints of the cusp in sheet $i$
are the classes at $i$ and $i+\eta(g)$, by
the actual transporter equation.  Every
character of $E_m$ evaluates that boundary
loop as the sum of its two endpoint values.
All boundary characters occur, so the
incidence map of this graph on $C_m$ has rank
at least $m-1$.  Its rank is $m-k$, where $k$
is the number of connected components.  Thus
$k=1$, and $\eta(g)$ generates $C_m$.
Varying $\eta$ proves primitivity at every
odd prime.

Suppose instead that $[g]\equiv0\pmod2$.
Pass to the characteristic finite class-two
quotient with $x^4=y^4=c^2=1$ and central
commutator.  Its normal forms are
$x^a y^b c^t$, $a,b\in\mathbb Z/4$, $t\in\mathbb F_2$.
The degree-two Frattini layer has basis
$X^2,Y^2,Z=[X,Y]$, and its induced action
depends only on $A_2\bmod2$.  For the
identity matrix, $1+\alpha$ has zero image.
For the nonidentity involution, choose an
integral orientation-preserving basis with
$X\mapsto X$, $Y\mapsto X+Y$ modulo degree
two.  The image of $1+\alpha$ is then
$\langle X^2+Z\rangle$.  Neither image
contains $Z$.  But the supposition puts
$g$ in this Frattini layer, and
$c=g\alpha(g)^{-1}$ would put $Z$ in that
image.  This contradiction proves dyadic
primitivity and therefore the first assertion.

In the two-class case, high-degree mod-two
cohomology of $E_m$ has dimension $m+2$:
the rank-two compact part and the rank-$(m-1)$
boundary part both have trivial mod-two
action.  Its $m$ cusp classes each contribute
one.  The other class splits into
$d=m/|\eta(B)|$ classes, all with centralizer
conjugate to $B\cap F_m$, giving
\[
 d\bigl(1+\dim H^1(B\cap F_m,\mathbb F_2)\bigr)=2.
\]
The integer $d$ is odd, so $d=1$.  Thus $B$
maps onto all odd invariant cyclic quotients.
At two, the nilpotent class $a/u$ in (C.8.1),
arising from the invariant integral character,
is nonzero on the other normalizer factor
and zero on the cusp factor.  Its restriction
to $B$ is nonzero.  The homology image of
$B$ on the dyadic fixed line therefore
contains a unit and is onto.  Its closed
image in $M^\alpha\cong\widehat{\mathbb Z}$
is onto modulo every prime, hence is the
entire group.  This proves the existence
of $a$ with $[a]=[g]$.

Complete $[g]$ to a basis $[g],e$ of $M$.
The involution matrix is
$\left(\begin{smallmatrix}1&b\\0&-1\end{smallmatrix}\right)$.
Killing $v$ kills $c=[g,v]$, so the remaining
quotient is abelian and equals
$M/(\alpha-I)M$.  In the two-class case
$b\in2\widehat{\mathbb Z}$ and its quotient
coordinates are
\[
 (r,t)\longmapsto(r+(b/2)t,\ t\bmod2)
      \in\widehat{\mathbb Z}\times C_2.
\]
The other section has nonzero second coordinate
by C.8 and zero first coordinate because it
is torsion.  Killing it leaves precisely
$\widehat{\mathbb Z}$, with $g\mapsto1$;
$B$ is onto by the homology-image assertion.
In the one-class case $b_2$ is odd, and
$(r,t)\mapsto2r+bt$ identifies the coinvariant
quotient with $\widehat{\mathbb Z}$ and sends
$g$ to $2$.  Prime by prime this map is
onto and its kernel is the image of
$\alpha-I$.  This proves all asserted
quotient comparisons.
\end{proof}

\statement{C.12}{Proved}{The exact pro-odd-kernel conclusions of LM4 D.78}
In the two-class situation of C.11, retain its actual
$B=C_F(w)$, $T=\langle w\rangle$, $\chi$ and chosen
$a\in B$.  Put $A=\overline{\langle a\rangle}$ and
$K=\ker(\chi|_B)$.  Then
\[
 \dim H^1(B_0,\mathbb F_2)=1\quad(B_0\leq B\text{ open}),
 \qquad B=K\rtimes A,
 \qquad A\cong\widehat{\mathbb Z},
                                                               \tag{C.12.1}
\]
and $K$ is pro-odd.  For every finite two-primary
$B$-module $Q$, inflation gives
\[
 H^n(\widehat{\mathbb Z},Q^K)\cong H^n(B,Q),                 \tag{C.12.2}
\]
and restriction to $A$ is the map induced by
$Q^K\hookrightarrow Q$, with $t$ acting on $Q$
as $a^t$.  The invariant-surjectivity criterion
is exactly C.3, and remains additional to these
conclusions.

\begin{proof}
Use C.6 for the one original cusp, and write
$\mathcal A_M,\mathcal B_M,\mathcal R_M$ for
its localized absolute, boundary and relative
cohomology.  C.8 and Brown restriction give
\[
 \mathcal A_M=\mathcal V_S(M)\oplus\mathcal U_T(M),\quad
 \mathcal V_S(M)=H^\bullet(S,M)[u^{-1}]_0,\quad
 \mathcal U_T(M)=H^\bullet(B\times T,M)[u^{-1}]_0.
\]
The actual boundary map is $(z,t)\mapsto(z,z)$,
with its second transport given by the original
$g$.  Let $e\in\mathcal A_k$ be the idempotent
selecting the $T$-factor.  It annihilates
$\mathcal B_M$, so the relative-to-absolute
map induces an isomorphism
$e\mathcal R_M\to e\mathcal A_M=\mathcal U_T(M)$.
The balanced perfect pairing of C.6 therefore
gives a perfect pairing on
$\mathcal U_T(M)\times\mathcal U_T(M^\vee)$,
with the actual cup product in $B\times T$.
Indeed its trace is the relative trace composed
with the inverse of that isomorphism.

For a finite $E$-set $X$, set $M=k[X]$ and
$Q_X=k[X^T]$.  The complementary basis summand
is free as a $k[T]$-module and is $B$-stable.
Its positive $T$-cohomology vanishes.  As
$\operatorname{cd}_2B\leq1$, it contributes
nothing in sufficiently high total degree.
The direct-product complex on the remaining
summand, on which $T$ acts trivially, gives
the natural isomorphism
\[
 \mathcal U_T(k[X])
       \cong H^0(B,Q_X)\oplus H^1(B,Q_X)u^{-1}.              \tag{C.12.3}
\]
The second summand is totally isotropic for
the perfect pairing, since its products lie
in $H^2(B,Q_X\otimes Q_X)=0$.  Consequently
$\dim H^1(B,Q_X)\leq\dim H^0(B,Q_X)$.
If the $B$-orbits in $X^T$ have stabilizers
$B_1,\ldots,B_s$, Shapiro identifies these
dimensions with $\sum_i\dim H^1(B_i,k)$
and $s$.  Every $\chi(B_i)$ is open in
$\widehat{\mathbb Z}$ and hence admits a
nonzero character to $C_2$.  Thus every
summand has dimension at least one, and
all have dimension exactly one.

Each open $B_0\leq B$ occurs as such a
stabilizer.  Choose $N\lhd E$ open with
$N\subseteq F$ and $N\cap B\subseteq B_0$.
In $X=E/(B_0TN)$ the identity coset is
$T$-fixed and its $B$-stabilizer is $B_0$.
Indeed quotient parity first removes the
$T$-factor, and then $N\cap B\subseteq B_0$
gives the equality.  This proves the
all-open-subgroup assertion in (C.12.1).

The powers of $a$ split $\chi$.  Its two-primary
subgroup $A_2\cong\mathbb Z_2$ lies in a Sylow
pro-two subgroup $P$ of $B$.  Two independent
characters $P\to C_2$ would factor through
$P/(P\cap N_0)$ for some open normal
$N_0\lhd B$ and extend to the open subgroup
$PN_0$, contradicting the rank just proved.
Thus every finite quotient of $P$ is cyclic,
by the Frattini criterion, and $P$ is procyclic.
Its map onto $\mathbb Z_2$ is an isomorphism;
since $A_2$ already maps onto, $P=A_2$.
All Sylow pro-two subgroups are conjugate,
and $\chi_2$ is injective on each.  Thus
$\ker\chi_2$, and its subgroup $K$, have
trivial Sylow pro-two subgroup and are
pro-odd.  Multiplication and its inverse
are explicitly
\[
 (k,t)\mapsto ka^t,\qquad
 b\mapsto(ba^{-\chi(b)},\chi(b)).
\]
Finally a pro-odd group has exact invariants
on finite two-primary coefficients, by
averaging over finite odd quotients, and
has no higher cohomology on them.  The
extension spectral sequence gives (C.12.2).
Restriction along the specified section
is precisely the coefficient inclusion
$Q^K\to Q$.  C.3 proves that replacing this
inclusion by an equality for all the
finite permutation coefficients is
equivalent to $K=1$; no such replacement
occurred in the argument.
\end{proof}

\statement{C.13}{Proved}{The exact transporter and conjugacy reduction of LM4 D.75}
For $\mu\in\widehat{\mathbb Z}^{\times}$ and
$i\in\{0,1\}$ put $v_i=c^iv$.  The unique solution
of $[a,v_i]=c^\mu$ is
\[
 a_{\mu,0}=c^{(\mu-1)/2}g,
 \qquad a_{\mu,1}=c^{(\mu+1)/2}g^{-1}.                        \tag{C.13.1}
\]
The full peripheral conjugacy to the corresponding
standard negative torus involution is equivalent to
\[
 \exists\mu\in\widehat{\mathbb Z}^{\times},\ i\in\{0,1\}:
                 c^{-\mu}a_{\mu,i}\sim_F a_{\mu,i}.           \tag{C.13.2}
\]
Existence in (C.13.2) is not asserted.

\begin{proof}
Every unit is odd at two, so the half-exponents
in (C.13.1) exist uniquely in
$\widehat{\mathbb Z}$.  Substituting
$gvg^{-1}=cv$ verifies both displayed
commutator equations.  Any two solutions
differ by the centralizer of $v_i$ in $F$,
which is trivial by C.8.  C.11 gives
$[a_{\mu,i}]=\pm[g]$, so C.5 applies to
each $a=a_{\mu,i}$.

Its image under $\alpha_i=\operatorname{Ad}(v_i)$
is $c^{-\mu}a$.  Thus (C.13.2) is equivalent
to the existence of $t\in C_E(a)\setminus F$.
Since $C_F(a)=\overline{\langle a\rangle}$,
this centralizer has two cosets and
$t^2=a^\kappa$ for a unique $\kappa$.
In the two-class case C.11 gives
$\chi(a)=\pm1$, hence $\kappa$ is twice
an element of $\widehat{\mathbb Z}$.
There is then a unique outside involution
$w=a^{-\kappa/2}t$ centralizing $a$.
It belongs to the other section class,
since the cusp class has trivial geometric
centralizer.  The standard two-class model
has generators $X,Y,v_0$ with
\[
 v_0Xv_0^{-1}=YXY^{-1},\qquad v_0Yv_0^{-1}=Y^{-1}.
\]
The homomorphism $X\mapsto a$, $Y\mapsto v_iw$,
$v_0\mapsto v_i$ preserves these relations
and sends $[X,Y]$ to $c^\mu$.

In the one-class case no outside involution
can centralize $a\ne1$, again by C.8.
Thus $\kappa_2$ is odd.  Replace $t$ by
$a^bt$ so that $\kappa+2b$ has its given
odd dyadic coordinate and coordinate one
at every odd prime.  This is a full unit
$\eta$.  The actual procyclic power
$z=(a^bt)^{\eta^{-1}}$ satisfies $z^2=a$
and lies outside $F$.  In this case the
standard model interchanges $X,Y$, and
$X\mapsto zv_i$, $Y\mapsto v_iz$,
$v_0\mapsto v_i$ preserves that action
and again sends $[X,Y]$ to $c^\mu$.

Both maps are full isomorphisms, by the
following direct boundary test.  If an
endomorphism of $\widehat F_2$ sends its
commutator to $c^\mu$ with $\mu$ a full
unit, its closed image contains $c$.
A proper image lies in a proper open
subgroup $U$.  The corresponding finite
punctured-torus cover has a degree-one
cusp represented by $c$.  If its degree
exceeds one there must be at least one
other cusp, and the class of $c$ is
nonzero primitive in $H_1(U,\widehat{\mathbb Z})$.
But $c^\mu$ is a commutator of elements
of $U$, a contradiction.  Thus the
endomorphism is onto and, by finite
generation and the finite-characteristic-
quotient Hopf argument, is an automorphism.
The extension maps are then isomorphisms
because they are identity on the quotient
label $C_2$.

Conversely a full peripheral isomorphism
of either standard model can be adjusted
by an inner automorphism so that its cusp
is $c^\mu$.  The image of its distinguished
cusp reflection is $c^t v$; conjugation
by a cusp power changes $t$ by twice an
arbitrary exponent, so it can be made
$v_i$ for $i=0$ or $1$.  Its standard
transporter is therefore the unique
$a_{\mu,i}$.  In the two-class model it
commutes with the other reflection; in
the one-class model it is the square of
an outside element.  Either supplies
an outside centralizer and hence (C.13.2).
This proves both directions with the
full peripheral labels retained.
\end{proof}

\statement{C.14}{Open}{The existence requested by prompt C.2}
For every negative involution of the full punctured-torus
loop-outer group, does (C.13.2) have a solution?
C.8 and C.13 prove uniqueness of every individual
commutator transporter and all claimed equivalences
in LM4 D.75.  They do not supply the additional
conjugacy.  An outside element in $C_E(a_{\mu,i})$
would solve it; absence of such elements for all
full units and both parity labels would refute it.
The equation remains a full profinite conjugacy
problem, and generation in each maximal pro-$p$
quotient does not solve it, as C.10 illustrates
for the related section-generation problem.

\statement{C.15}{Proved}{Prime-order lifting and full closed-surface primary faithfulness}
Every outer action of $C_p$ on the full free profinite group
$\widehat F_2$ has an actual order-$p$ representative if its
outer image has order $p$.  For the full completed closed
surface group $\Pi=\widehat{\pi_1(S_g)}$, $g\geq2$, the map
\[
 \operatorname{Aut}(\Pi)\longrightarrow
       \operatorname{GL}(H_1(\Pi,\mathbb Z_p))
                                                               \tag{C.15.1}
\]
is injective on finite $p$-subgroups.  This second assertion
concerns actual automorphisms, not outer automorphisms.

\begin{proof}
The centre-free extension attached to the specified rank-two
outer action has an open subgroup $\widehat F_2$ of index
$p$.  If this extension were torsion-free, Theorem~1.2 of
\cite{CEJarden1982} would force $p$ to divide $2-1$.
Thus it has torsion.  A finite cyclic subgroup injects into
its quotient $C_p$, since $\widehat F_2$ is torsion-free;
it consequently has order $p$ and is a section.  Conjugation
by its labelled generator is the required genuine lift.
This argument uses the full profinite rank-divisibility
result, not a pro-$p$ analogue of the desired conclusion.

Put $P=\Pi^{(p)}$.  The closed lower-central quotients
$P/\gamma_nP$ are torsion-free and $\bigcap_n\gamma_nP=1$;
these are Propositions~2.14--2.15 of
\cite{CEMorales2024} for the orientable surface relation.
For $n\geq1$, put
$J_n=\ker(\operatorname{Aut}(P)\to
\operatorname{Aut}(P/\gamma_{n+1}P))$.  The map
\[
 b\longmapsto\bigl(z\longmapsto b(z)z^{-1}
                         \bmod\gamma_{n+2}P\bigr)
\]
embeds $J_n/J_{n+1}$ in the torsion-free group
$\operatorname{Hom}_{\rm cont}(P^{\rm ab},
\gamma_{n+1}P/\gamma_{n+2}P)$.  Centrality of the target
makes it additive in $z$.  An IA automorphism acts trivially
on each lower-central graded piece, which makes it additive
in $b$.  Its kernel is precisely $J_{n+1}$.
A finite-order element of $J_1$ therefore lies in every
$J_n$ and is trivial.

Suppose now that $a\in\operatorname{Aut}(\Pi)$ has order
$p$ and acts trivially on $H_1(\Pi,\mathbb Z_p)$.
The preceding calculation makes its induced action on $P$
trivial.  There is a specified continuous quotient
\[
 E=\Pi\rtimes_a C_p\longrightarrow P\times C_p,
                     \qquad(h,s^j)\longmapsto(\bar h,s^j).
                                                               \tag{C.15.2}
\]
The quotient $\Pi\to P$ is a mod-$p$ cohomology-ring
isomorphism.  In degrees zero and one this is immediate;
in degree two it follows from the nondegenerate surface
cup pairing on the common $2g$ generator characters;
higher cohomology vanishes by the full and pro-$p$
surface duality comparisons.  Hochschild--Serre comparison
for (C.15.2) thus gives
\[
 H^*(E,\mathbb F_p)\cong H^*(P\times C_p,\mathbb F_p),
 \qquad \dim H^n(E,\mathbb F_p)=2g+2\quad(n>2).
                                                               \tag{C.15.3}
\]
For clarity, the surface comparison can be obtained from
finite-cover cohomological effacement: trivialize a finite
module by a cover, kill degree-one classes by their finite
character kernels, and kill the orientation class by a
cover of degree divisible by its coefficient exponent.
This proves goodness and the full surface cohomology used
here; the pro-$p$ presentation has the same one-relation
resolution and the same cup matrix.

Let $\eta$ be the surface class and $u$ the degree-two
periodicity class of $C_p$.  The class $\eta u^m$ is
nonzero in (C.15.3).  A subgroup $S$ of order $p$ in $E$
projects isomorphically to $C_p$ and maps to that factor
in (C.15.2), since $P$ is torsion-free.  Its normalizer is
$C_\Pi(S)\times S$: projection to the abelian quotient
makes every normalizer element centralize $S$.
Theorem~2 of \cite{CEKochloukovaZalesskii2007} applies to
$E$, since its virtual $p$-cohomological dimension is two
and it contains no $C_p^2$.  Its dense restriction
embedding in degrees greater than two is an isomorphism
here, because its source is finite.  Nonvanishing of
$\eta u^m$ therefore gives a factor $C=C_\Pi(S)$ with
$\eta|_C\ne0$.  Nondegeneracy of the surface pairing
makes $H^1(\Pi,\mathbb F_p)\to H^1(C,\mathbb F_p)$
injective.  This one normalizer already contributes at
least $1+2g+1$ dimensions in every degree greater than
two, while (C.15.3) gives exactly $2g+2$.  Consequently
there is one section class, and
\[
 \dim H^1(C,\mathbb F_p)=2g,\qquad H^2(C,\mathbb F_p)\ne0.
                                                               \tag{C.15.4}
\]
The assertion applies to the specified section as well:
conjugation under $E$ and under $\Pi$ gives the same
section classes.

We verify directly that (C.15.4) forces $C=\Pi$.
Continuous cohomology satisfies
\[
 H^j(C,\mathbb F_p)=
 \varinjlim_{C\leq U\leq_o\Pi}H^j(U,\mathbb F_p),
                                                               \tag{C.15.5}
\]
with restriction maps.  Indeed a finite-valued continuous
cochain on $C$ factors through $C/(C\cap V)$ for some
open normal $V\leq\Pi$ and extends to $CV$ through
$CV/V\cong C/(C\cap V)$.  Choosing a common $V$ for
cochains and their coboundaries proves (C.15.5).
Restriction on second cohomology between these covering
surface groups multiplies by the covering index.  A
nonzero class in the direct limit has a representative
at $U$ such that every smaller $W\geq C$ has
$p\nmid[U:W]$.  Restriction on first cohomology is then
injective by transfer, including into the direct limit.
If $C$ had infinite index, the quantities
\[
 \dim H^1(W,\mathbb F_p)=2+2[\Pi:W](g-1)
\]
would be unbounded inside the fixed $2g$-dimensional
space in (C.15.4).  Thus $C$ is open, and the same
formula gives $[\Pi:C]=1$.  Hence $a=1$.
A nontrivial kernel on a finite $p$-group contains an
element of order $p$, so (C.15.1) is injective as stated.
\end{proof}

\statement{C.16}{Proved}{The order-three assertion of LM4 D.42}
Every element of order three in the full three-cusp
loop-outer group allowing cusp permutations has
multiplier $1\in\widehat{\mathbb Z}^{\times}$.

\begin{proof}
By the refereed pure-group assertion R1, its cusp
permutation is a three-cycle.  Label the boundary
classes $x,y,z=(xy)^{-1}$ so that they go successively
to $y,z,x$, with common power $\lambda$.  C.15 supplies
an actual order-three lift $\alpha$ of this outer action.
The equality $\lambda^3=1$ forces $\lambda_3=1$,
since $\mathbb Z_3^{\times}$ has no element of order
three.  If $\lambda_q\ne1$ at another prime, its
reduction $\zeta\in\mathbb F_q^{\times}$ has order
three.  Thus $q\equiv1\pmod3$ and $q\geq7$.
The continuous epimorphism
\[
 \chi:F\longrightarrow C_q,\qquad
 \chi(x)=\zeta,\quad\chi(y)=1,\quad\chi(z)=\zeta^2
                                                               \tag{C.16.1}
\]
satisfies $\chi\alpha=\chi$, by
$1+\zeta+\zeta^2=0$ and direct substitution of the
three powered boundary classes.  Set $N=\ker\chi$.
Its corresponding degree-$q$ cover has one cusp above
each original cusp and genus $h=(q-1)/2$.
Filling those three cusps is the exact full profinite
quotient
\[
 \Pi_h=N/
 \overline{\langle\!\langle x^q,y^q,z^q\rangle\!\rangle}_N
       \cong\widehat{\pi_1(S_h)}.
                                                               \tag{C.16.2}
\]
This is completion of the finite-cover presentation;
completion is right exact on the indicated quotient.
The action induces an actual automorphism $b$ of order
dividing three on $\Pi_h$.  The deck generator induced
by conjugation with $y$ has order $q$ after filling.
Because $\alpha$ induces the identity on $F/N$, its
commutator with every deck map is inner.  Accordingly
$b_*$ commutes with the deck representation on homology.

The regular cover of the two-loop rose gives the rational
deck-module class $[H_1(N,\mathbb Q_3)]=[1]+[\mathbb Q_3[C_q]]$.
Its boundary submodule is trivial of dimension two,
because each of its three cusps has full inertia $C_q$
and their sole relation is their sum.  The filling exact
sequence therefore identifies
\[
 H_1(\Pi_h,\mathbb Q_3)
       \cong\mathbb Q_3[C_q]/\mathbb Q_3
       \cong\mathbb Q_3[X]/(1+X+\cdots+X^{q-1}).
                                                               \tag{C.16.3}
\]
Semisimplicity justifies passage from the computed
representation class to this module isomorphism.  It
is cyclic, so its deck-commuting endomorphism algebra
is the algebra on the right.  Each field factor is
unramified over $\mathbb Q_3$, because it is generated
by roots of unity of order prime to three.  An
unramified extension of $\mathbb Q_3$ contains no
nonidentity cube root of unity: its residue multiplicative
group has order $3^f-1$, and the principal units have
no three-torsion by the valuation of $(1+3^rv)^3-1$.
Hence the order-three unit $b_*$ in this algebra is
one.  It is also identity on integral $\mathbb Z_3$
homology, which embeds in rational homology.

C.15, now for the full closed surface group and $p=3$,
gives $b=1$.  Its action on the orientation class,
however, is $\lambda$.  To verify this last comparison,
use the relative fundamental class before filling:
its boundary is the sum of the three oriented cusp
generators.  The action permutes them and raises all
to the common power $\lambda$, so it multiplies the
relative class by $\lambda$.  Filling identifies this
class with the closed orientation class.  This
contradicts $\lambda_q\ne1$ and proves the assertion.
\end{proof}

\statement{C.17}{Proved}{The primary fixed-group and cusp arguments needed for the quartic audit}
Let $a$ be an actual automorphism of order $p$ of a full
free profinite group $F_r$ which induces the identity on
$F_r^{(p)}$.  Then $F_r\rtimes_a C_p$ has one
$F_r$-conjugacy class of sections, and restriction gives
\[
 H^1(F_r,\mathbb F_p)\xrightarrow{\sim}
 H^1(C_{F_r}(a),\mathbb F_p).                              \tag{C.17.1}
\]
If $F_r$ is the group of a sphere with $r+1\geq3$ cusps
and $a$ preserves each cusp class with $p$-adic multiplier
one, then $a=1$.  If $r=2$ and $a$ preserves the torus
cusp, the same conclusion holds.  In particular the
homology representation of the full torus loop-outer
group is injective on each finite $p$-subgroup at its
$p$-adic coordinate.

\begin{proof}
The quotient $F_r\to F_r^{(p)}$ is a mod-$p$ cohomology
isomorphism: both groups have cohomological dimension
one and the same $r$ generator characters.  The specified
map $F_r\rtimes_a C_p\to F_r^{(p)}\times C_p$ therefore
induces a cohomology-ring isomorphism by
Hochschild--Serre comparison.  In every degree greater
than one its dimension is $r+1$.  Multiplication by the
periodicity class $u\in H^2(C_p,\mathbb F_p)$ embeds the
$r$ generator characters into high cohomology.  These
classes vanish on each section, since its image in the
torsion-free pro-$p$ factor is trivial.
Brown restriction, with normalizers $C_{F_r}(S)\times S$,
therefore gives
\[
 k+\sum_{[S]}\dim H^1(C_{F_r}(S),\mathbb F_p)=r+1,
 \qquad \sum_{[S]}\dim H^1(C_{F_r}(S),\mathbb F_p)\geq r.
\]
Here $k\geq1$ is the number of section classes; the
restriction theorem is applicable because $F_r$ is open
of dimension one and finite subgroups inject into $C_p$.
Thus $k=1$ and (C.17.1) follows, including the claim that
it is the actual restriction map.

We record the needed pro-$p$ elementary facts.  A free
pro-$p$ group embeds in the units of the completed
noncommutative power-series algebra by sending its free
generators to $1+X_i$.  The universal finite nilpotent
quotients, or equivalently Hall collection of commutators,
show that its lower-central graded pieces are torsion-free
$\mathbb Z_p$-modules and separate its elements.  The
successive-deviation maps used in C.15 therefore show
that its IA automorphism kernel is torsion-free.  The
centralizer of a nontrivial element is procyclic: a closed
subgroup is free pro-$p$, and a noncyclic free pro-$p$
group has trivial centre, as follows by commuting its
Magnus series with two distinct free generators.  For a
sphere cusp this procyclic centralizer is the cusp itself,
since the cusp is primitive.  The torus commutator is
also not a proper $p$th power.  Its first nonzero mod-$p$
Magnus term is $XY-YX$, of degree two.  A $p$th power
has initial degree divisible by $p$; for $p=2$, the
only possible degree-two square $(aX+bY)^2$ has square
terms $a^2X^2+b^2Y^2$ which prevent it from being
$XY+YX$.  Thus its centralizer is its whole cusp as well.

For either surface structure, the full peripheral
normalizer is an extension of its full procyclic cusp
by $C_p$.  For a lift $v$ write
$v c v^{-1}=c^\lambda$ and $v^p=c^t$.
Its image in $F_r^{(p)}\times C_p$ has first coordinate
$c^w$ by the preceding centralizer calculation; hence
$t_p=pw$.  Where $\lambda_q\ne1$, the identities
$\lambda^p=1$ and $(\lambda-1)t=0$ give $t_q=0$ and
$1+\lambda_q+\cdots+\lambda_q^{p-1}=0$.  Where
$\lambda_q=1$, that norm is $p$, and $t_q$ is divisible
by $p$, including at $q=p$ by the specified comparison.
A coordinatewise norm correction thus makes $vc^z$ a
section centralizing the entire cusp $\mathbb Z_p$.
By uniqueness of its class, a conjugate of that cusp
$\mathbb Z_p$ lies in $C=C_{F_r}(a)$.

Here is the closed-subgroup dimension argument with the
covering maps retained.  If $C$ is open, Schreier gives
$\dim H^1(C,\mathbb F_p)=1+[F_r:C](r-1)$, so (C.17.1)
forces index one.  If $C$ has infinite index but finite
$p$-part of its index, choose an open neighbourhood of
$C$ where that part is maximal.  All further restriction
maps are injective by prime-to-$p$ transfer; the
restriction-colimit argument (C.15.5) then embeds
unbounded open-subgroup first-cohomology dimensions in
$H^1(C,\mathbb F_p)$, a contradiction.  In the remaining
case choose $U\geq C$ open with $p\mid[F_r:U]$.
A selected full cusp $\mathbb Z_p$ contained in $C$
gives a boundary of this finite cover of degree prime
to $p$.  Its class cannot be the unique boundary above
its original cusp, since the sum of the boundary
degrees there is $[F_r:U]$.

In the torus case there are consequently at least two
boundaries, and the selected boundary has nonzero
primitive homology.  A character detecting it restricts
to $C$ and detects the selected cusp $\mathbb Z_p$.
This contradicts (C.17.1), since every character of
$F_2$ vanishes on its commutator cusp.  More explicitly,
if that boundary has degree $d$ prime to $p$, its
parameter on the contained cusp $\mathbb Z_p$ is
multiplication by $d^{-1}$, a nonzero scalar modulo $p$.
In the sphere case, apply the construction to all
$r+1$ distinct original cusps.  There are at least
$r+2$ boundaries in $U$, so the $r+1$ selected boundary
classes are independent: the only relation among all
oriented boundary classes is their sum.  Their characters
restrict independently to $C$, contradicting dimension
$r$ in (C.17.1).  Thus in both cases $C=F_r$.

For an outer torus element of order $p$ acting trivially
on $p$-adic homology, C.15 gives an actual order-$p$
lift.  Its pro-$p$ action is identity by torsion-freeness
of the IA kernel just proved, and the torus argument
makes the lift identity.  A nontrivial kernel on a
finite $p$-subgroup contains an order-$p$ element,
which proves the final assertion.
\end{proof}

\statement{C.18}{Proved}{The global torus involution sign, with the Heisenberg comparison specified}
Every involution of the full torus loop-outer group has
multiplier equal to its global dyadic sign, $1$ or $-1$.
This proves the upstream sign assertion LM4 D.43 used
in the negative-involution and quartic calculations.

\begin{proof}
Take an actual order-two representative $\alpha$ by C.15.
The full class-two Heisenberg quotients identify its
peripheral multiplier with the determinant of its
abelianization matrix, so $\lambda^2=1$.
For an odd prime $q$, use the characteristic quotient
\[
 G_q=\mathbb F_q^2\times\mathbb F_q,\qquad
 (v,z)(w,t)=(v+w,z+t+\tfrac12\omega(v,w)),
\]
where $\omega((a,b),(a',b'))=ab'-ba'$, with $x,y$
the standard horizontal generators and $c$ the central
coordinate one.  This is exactly the class-two,
exponent-$q$ quotient: collection gives at most $q^3$
normal forms, and these coordinates give all $q^3$.
Its kernel is therefore characteristic.
The associated degree-$q^3$ punctured-torus cover has
$q^2$ cusps of degree $q$.  Fill all of them.  The
resulting full closed surface group has genus
$g_q=1+(q^3-q^2)/2$, and $\alpha$ induces an actual
involution on it.  The quotient map is the closed
normal quotient by all $G_q/\langle c\rangle$-labelled
conjugates of $c^q$; completion of the finite-cover
presentation proves that this is the full surface
completion, not merely its pro-$p$ image.

Let $V=H_1(\Pi_q,\mathbb Q_p)$ and $A=\alpha_*$.  The
deck representation $\rho$ and its induced automorphism
$\theta$ satisfy
\[
 A\rho(h)A^{-1}=\rho(\theta(h)),\quad
 A^2=1,\quad\theta(c)=c^{\lambda\bmod q}.
                                                               \tag{C.18.1}
\]
The graph-cover chain complex and the exact sequence for
filling the cusps give the actual representation class
\[
 [V]=2[1]+[\mathbb Q_p[G_q]]
                 -[\mathbb Q_p[G_q/\langle c\rangle]].
                                                               \tag{C.18.2}
\]
The surface alternating form is deck-invariant and has
multiplier $\lambda_p$ under $A$.  To check the latter
without realizing $\alpha$ geometrically, note that its
permutation of the odd number $q^2$ of cusps has a fixed
cusp.  Fill the others first.  The remaining boundary
word of the genus-$g_q$ one-cusp group has nondegenerate
degree-two commutator bivector.  Raising that boundary
to power $\lambda$ multiplies this bivector, and hence
the inverse alternating form, by $\lambda_p$.
Filling the last cusp preserves first homology.

Over a characteristic-zero splitting field, the
$c$-eigenspace $V_j$ with eigenvalue $\zeta_q^j$,
$j\ne0$, has dimension $q^2$, by (C.18.2).  It pairs
nondegenerately only with $V_{-j}$.  If $\lambda_q=-1$,
then $A$ exchanges these spaces.  Write $B,C$ for the
two exchange matrices and $P$ for their pairing matrix.
They satisfy $BC=CB=1$ and
$-B^{\mathsf T}P^{\mathsf T}C=\lambda_pP$.
Determinants in the odd dimension $q^2$ give
$-1=\lambda_p^{q^2}$, hence $\lambda_p=-1$ for every
$p$.  Thus any negative odd coordinate forces the
global negative sign.

It remains to exclude a negative dyadic coordinate
with all odd coordinates positive.  Assume
$\lambda_3=1$.  On $G_3/\langle c\rangle$ the
involution has determinant one and is $I$ or $-I$.
An automorphism of $G_3$ with this action has form
$\theta(v,z)=(Mv,z+\ell(v))$.  If $M=I$, order two
gives $\ell=0$.  If $M=-I$, conjugation by an inner
automorphism from $(u,0)$ changes $\ell$ by
$-2\omega(u,\cdot)$, so choose $u$ to make it zero.
Lift $u$ to $F$ and conjugate the actual automorphism
$\alpha$ by that inner automorphism; its order remains
two.  Thus $\theta$ is $1$ or $(v,z)\mapsto(-v,z)$.

Over $L=\mathbb Q_2(\zeta_3)$, each nontrivial central
block of $G_3$ is $M_3(L)$.  Explicitly, on basis
$e_r$, $r\in\mathbb F_3$, its matrices are
\[
 Xe_r=e_{r+1},\qquad Ye_r=\zeta_3^{-jr}e_r,
                    \qquad Ze_r=\zeta_3^je_r.
\]
The distinct $Y$-eigenprojectors and powers of $X$
span the full matrix algebra, proving the block
assertion.  By (C.18.2), $V_j=W_j\otimes U_j$ with
$\dim U_j=3$.  The rational isometry $J=I$, or
$Je_r=e_{-r}$, implements the two possible $\theta$.
Consequently $A|_{V_j}=J\otimes B_j$ with $B_j^2=I$.
Choosing conjugate bases of the two multiplicity spaces
gives $B_2=\sigma(B_1)$, where
$\sigma(\zeta_3)=\zeta_3^{-1}$.  Their invariant
cross-pairing is $I_3\otimes S$ for an invertible $S$,
so
\[
 B_1^{\mathsf T}S\sigma(B_1)=\lambda_2 S.
\]
Since $\det B_1\in\{1,-1\}$ is fixed by $\sigma$,
determinants give $1=\lambda_2^3$, and hence
$\lambda_2=1$.  Together with the negative-coordinate
argument this proves the asserted global sign.
\end{proof}

\statement{C.19}{Proved}{Positive torus sections, full primary square, and actual quartic lifting}
A positive torus involution acts as $-I$ on
$H_1(F,\mathbb Z_p)$ for every prime $p$.  Its
centre-free full extension has three kernel-conjugacy
classes of sections, each with
$H^1(C_F(S),\mathbb F_2)=0$.  Its cusp normalizer is
torsion-free, and a labelled cusp lift $v$ satisfies
$v^2=c^\kappa$ with $\kappa_2$ a unit.
Every order-four outer torus action has an actual
order-four lift.  For such an action,
$A_p^2=-I$ at every prime, $\lambda^2=1$, and
$\lambda_2=1$.  These are precisely the portions of
LM4 D.46 and D.55 needed below.

\begin{proof}
For the positive involution, C.17 excludes the identity
matrix on dyadic homology.  A nonidentity involution
in $\operatorname{GL}_2(\mathbb Q_2)$ with determinant
one is $-I$, so $A_2=-I$.  Lift the outer action by
C.15 and write $E_b=F\rtimes\langle s\rangle$.
Its two mod-two generator characters $\chi_1,\chi_2$
extend by value zero on $s$, and its quotient character
is $u$.  The two-row split spectral sequence gives
basis $u^n,u^{n-1}\chi_1,u^{n-1}\chi_2$ in every
positive degree.  For an integral lift
$\varphi_i:F\to\mathbb Z/4$ one has
$\varphi_i\beta=-\varphi_i$.  The cochain
$\widetilde\chi_i(hs^\epsilon)=\varphi_i(h)$ has
coboundary $2\epsilon\varphi_i(k)$ on
$(hs^\epsilon,ks^\eta)$.  Its Bockstein is therefore
$u\chi_i$; the Bockstein of a degree-one mod-two
character is its cup square.  Thus the localized
three-dimensional algebra is generated by the
idempotents $\chi_i/u$ and is reduced.

Brown restriction identifies that algebra with
\[
 \prod_{[S]}\bigl(\mathbb F_2\oplus
 H^1(C_F(S),\mathbb F_2)u^{-1}\bigr).
                                                               \tag{C.19.1}
\]
Each second summand has square zero because its group
is a closed subgroup of the dimension-one group $F$.
Reducedness forces all these summands to vanish;
dimension three gives exactly three section classes.
Their evaluations on the two extended characters are
three distinct points of $\mathbb F_2^2$.

No closed $C\leq F$ with $H^1(C,\mathbb F_2)=0$
contains the full cusp $\mathbb Z_2$.  Indeed finite
2-part of its index contradicts the restriction-colimit
and prime-to-two transfer argument of C.17.  Infinite
2-part supplies a cover of even degree containing that
cusp $\mathbb Z_2$; its selected boundary has odd
degree, there is another boundary, and its primitive
boundary character restricts nontrivially to $C$.
By C.18 the multiplier is globally one.  A section
in the cusp normalizer would therefore centralize the
full cusp, contradicting (C.19.1).  A labelled cusp
lift has $v^2=c^\kappa$; if $\kappa_2$ were even,
solving $2z=-\kappa$ in $\widehat{\mathbb Z}$ would
make $vc^z$ a section.  Hence $\kappa_2$ is odd.

Suppose a nonzero invariant character
$\chi:F\to C_q$ existed at an odd prime.  Extend it
by value zero on $s$.  The identity $v^2=c^\kappa$
gives $2\widetilde\chi(v)=0$, so $v$ lies in
$E_N=N\rtimes\langle s\rangle$, $N=\ker\chi$.
The degree-$q$ cover has type $(1,q)$.  Over
$\mathbb Z_2$, deck averaging splits its homology
into a rank-two handle summand, mapping to base
homology by the odd-index inclusion, and the boundary
augmentation summand of rank $q-1$.  The action is
$-I$ on the first summand.  It is identity on the
second: the cusp-normalizing $v\in E_N$ fixes each
cusp label and its oriented boundary class.  Thus
$I+\beta=0$ modulo two on the entire homology.
Writing $v=ns$ gives $[v^2]=[n\beta(n)]=0$ modulo
two, whereas $[c^\kappa]\ne0$ because $c$ is a
primitive boundary and $\kappa_2$ is odd.  This
contradiction excludes the invariant character.
At every odd prime the matrix is an involution with
determinant one, hence $I$ or $-I$.  The first choice
would supply such a character, so $A_p=-I$ everywhere.

Now form the centre-free extension $E$ of an outer
$C_4$ action, and let $E_2$ be the preimage of its
square.  The square is dyadically positive, since
$\mathbb Z_2^\times$ has no element of order four.
It therefore has the three section classes just
computed.  The quotient $E/E_2=C_2$ fixes at least
one of these three classes.  Choose a representative
$S$ and adjust a lift of a generator by an element
of $F$ to normalize $S$.  Put $C=C_F(S)$ and
$J=N_E(S)/S$.  The exact sequences are
\[
 1\to C\to N_E(S)\to C_4\to1,
 \qquad 1\to C\to J\to C_2\to1,
 \qquad S\subset Z(N_E(S)).
                                                               \tag{C.19.2}
\]
The first-cohomology vanishing gives $C^{(2)}=1$,
so $J^{(2)}=C_2$.  If $J$ were torsion-free, its
open projective subgroup $C$ would make $J$
projective by Lemma~2.1 of \cite{CEJarden1982}.
Its maximal pro-two quotient would then be free
pro-two, contradicting $J^{(2)}=C_2$.
A torsion element in $J$ must have order two and
lie outside $C$.  Its lift $n$ in $N_E(S)$ has
$n^2\in S$ and image of order four in $C_4$,
so $n^2$ is the nonidentity element of $S$ and
$n$ has order four.  Inverting it if necessary
produces the labelled section and actual lift.
The square now acts as $-I$ at every prime by the
positive-involution result.  Hence $A_p^2=-I$.
The global sign theorem for the square gives
$\lambda^2=1$; over $\mathbb Q_2$, $A_2$ has
minimal polynomial $X^2+1$, giving determinant
$\lambda_2=1$.  This completes all assertions.
\end{proof}
